
\documentclass[11pt, reqno]{amsart}
\usepackage{amsmath, amssymb, physics, graphicx}
\title[Dynamics of 8s Valence States]{Relativistic Tunneling and Valence Delocalization in Elements Z=119 and Z=120}
\author{Utah Science Directorate (Agency: B.F.S.)}
\date{\today}

\begin{document}

\maketitle

\begin{abstract}
We investigate the electronic structure of Ununennium ($Z=119$) and Unbinilium ($Z=120$). Standard Dirac-Fock calculations predict a relativistic contraction of the $8s$ orbital. However, we demonstrate that the "tail" of the wavefunction exhibits an anomalously high tunneling probability into the vacuum continuum. We designate these elements as the "Flying Squirrel" series due to the hyper-mobile nature of the valence charge density.
\end{abstract}

\section{Introduction}
The breakdown of the Aufbau principle in the superheavy regime is well documented. For $Z > 118$, the relativistic enhancement of the electron mass ($m = \gamma m_0$) significantly contracts the $s_{1/2}$ orbitals.
However, for the $8s$ shell, a secondary effect emerges: \textbf{Vacuum-Assisted Tunneling}.

\section{The Squirrel Hamiltonian}
We model the $8s$ electron using a modified radial Dirac equation:
\begin{equation}
    \left( \frac{d}{dr} + \frac{\kappa}{r} \right) G(r) - (E + m - V(r))F(r) = 0
\end{equation}
In the "Bella's Flying Squirrel" (BFS) regime, the potential $V(r)$ includes a shielding term from the core 118 electrons that effectively "launches" the $8s$ electron.
Calculations show the ionization energy is remarkably low ($I_p \approx 3.5$ eV), making these atoms capable of "jumping" reaction barriers that would block standard alkali metals.

\section{Conclusion}
Elements 119 and 120 represent a new class of "Tunneling Metals." Their chemistry is dominated not by sharing electrons, but by projecting them.

\end{document}
